% Part: computability
% Chapter: tm-computations
% Section: church-turing-thesis

\documentclass[../../../include/open-logic-section]{subfiles}

\begin{document}

\olfileid{tur}{mac}{ctt}
\olsection{The Church--Turing Thesis}

图灵机旨在成为能行过程这一概念的精确替代。图灵认为，任何同时理解了能行过程概念和图灵机概念的人，都会产生这样的直觉：任何能通过能行过程完成的事情，也都能由图灵机完成。这一主张得到了以下事实的支持：所有其他被提出用于替代能行过程概念的精确方案，最终都被证明与图灵机概念外延等价——也就是说，它们恰好能计算相同的函数集。这一主张被称为\emph{邱奇--图灵论题}。

\begin{defn}[Church--Turing 论题]
\emph{邱奇--图灵 论题}断言：任何能行可计算的事物都是图灵可计算的。
\end{defn}

对 邱奇--图灵 论题的援引有两种方式。第一种用法是作为一种偷懒的借口。假设我们有一个计算某事物的能行过程的描述，比如用“伪代码”写的。那么即使我们实际上并未构造出该图灵机，也可以援引 邱奇--图灵 论题来证成如下主张：存在某台图灵机计算了相同的函数。

邱奇--图灵 论题的另一种用法在哲学上更为有趣。可以证明，存在一些图灵机无法计算的函数。据此，利用 邱奇--图灵 论题，就可以得出结论：它无法通过任何能行过程来计算。因为如果存在这样的能行过程，根据 邱奇--图灵 论题，就会推出存在一台计算它的图灵机。因此，如果我们能证明没有图灵机可以计算它，那么也就不可能存在计算它的能行过程。具体而言，援引 邱奇--图灵 论题可以主张：所谓的停机问题不仅无法由图灵机解决，而且根本无法通过能行过程解决。

\end{document}